\noindent
\graphicspath{{./16-07/Images/}} 
%%%%%%%%%%%%%%%%
% Page de gauche
%%%%%%%%%%%%%%%%
\begin{minipage}[t][\textheight][t]{0.33\textwidth}
    \colorbox{StrongBlue}{
        \parbox[c][\textheight][c]{\textwidth}{ 
            \color{OffWhite} 
            \begin{center}
                \begin{minipage}[t][0.6\textheight][c]{0.8\textwidth}
                    {\huge \textbf{Dimanche 17 Août}}\\
                    \normalsize
                    \begin{center}
                        \begin{tabular}{@{} l @{\hspace{0.5cm}} l @{}}
                            Abri de départ : &\hfill Crouesty \\ 
                            Abri d'arrivée : &\hfill  Hoëdic \\ 
                            Distance : &\hfill 13 Milles nautiques \\ 
                            Départ : &\hfill 14 h 00 \\ 
                            Arrivée prévue : &\hfill 18 h 00 \\ 
                            Lever soleil 
\includegraphics[width=8px,height=8px]{../../Pictos/soleil_OffWhite.png} : &\hfill 07 h 11 \\ 
                            Coucher soleil 
\includegraphics[width=8px,height=8px]{../../Pictos/lune_OffWhite.png} : &\hfill 21 h 18 \\ 
                            \\
                            Abris replis : &  \\
                                    & • Port du Crouesty \\   
                                    & • Port de saint Goustan \\
                                    & • Port de l'Argol \\
                                            
                                
                        \end{tabular}
                    \end{center}
                    
                \end{minipage}%

                \begin{minipage}[t][0.4\textheight][t]{\textwidth}

                    \begin{center}
    {\large \textit{Marées \textbf{Crouesty}}}\\ \vspace{0.5em}
    
    \renewcommand{\arraystretch}{1.40} % Adjust the row height by a factor of 1.4
    
    % \begin{tabular}{ |C{5em}|c|C{4em}|C{4em}| } \hline
    %     & Coef & Heure & Hauteur\\ \hline
    %     BM & \multirow{2}{1em}{43}  & 05 h 07 & 2.09 m\\ \cline{1-1} \cline{3-4}
    %     PM &                        & 12 h 05 & 4.21 m\\ \hline
    %     BM & \multirow{1}{1em}{43}  & 17 h 51 & 1.85 m\\ \hline
    % \end{tabular}

    \begin{tabular}{ |C{5em}|c|C{4em}|C{4em}| } \hline
        & Coef & Heure & Hauteur\\ \hline
        BM & \multirow{2}{1em}{49}  & 05 h 28 & 1.79 m\\ \cline{1-1} \cline{3-4}
        PM &                        & 11 h 40 & 4.26 m\\ \hline
        BM & \multirow{1}{1em}{49}  & 18 h 11 & 1.92 m\\ \hline
    \end{tabular}
\end{center}

                    \begin{center}
    {\large \textit{Marées \textbf{Hoëdic}}}\\ \vspace{0.5em}
    
    \renewcommand{\arraystretch}{1.40} % Adjust the row height by a factor of 1.4
    \begin{tabular}{ |C{5em}|c|C{4em}|C{4em}| } \hline
        & Coef & Heure & Hauteur\\ \hline
        BM & \multirow{2}{1em}{49}  & 05 h 20 & 1.66 m\\ \cline{1-1} \cline{3-4}
        PM &                        & 11 h 34 & 4.03 m\\ \hline
        BM & \multirow{1}{1em}{49}  & 18 h 02 & 1.75 m\\ \hline
    \end{tabular}
\end{center}

                \end{minipage}%
            \end{center}       
        }      
    }
\end{minipage}%
\begin{minipage}[c][\textheight][t]{0.10\textwidth}
\hfill
\end{minipage}%
\begin{minipage}[c][\textheight][t]{0.50\textwidth}
    \vspace{2cm}
    \color{StrongBlue}    
    {\huge \textbf{Description de la navigation}}\\
            \begin{itemize}
                % \item Journée de navigation relativement courte et concentrée sur l'après-midi.  
                \item Un des objectifs du camp est de permettre aux jeunes qui le souhaitent de vivre au moins une messe. L'église d'Arzon située à 20 minutes du port a été contactée et il est convenu que les jeunes puissent animer une partie de l'office prévu à 10h30. Retour sur les bateaux prévu pour midi suivi du repas. Un.e chef.fe réaprovisionne le frais pendant la matinée au super U d'Arzon. 
                \item Départ du port à 14h, hissage une fois sortit du chenal entre 14h30 et 14h45. Cap plein sud d'un seul bord si le vent le permet.
                \item Après 30 minutes de navigation lancement à la VHF du jeu "escape boat" qui a été mis en place le matin pendant que les jeunes étaient à la messe.  
                \item Une fois que la pointe de Houat est dépassée pleine Ouest et largement avant la cardinale Est  \textit{Men Groise}, prendre un cap Sud Est en direction d'Hoëdic. 
                \item Veiller à passer entre la marque spéciale EDF 4 et le danger isolé de \textit{La chèvre} lors de l'approche de l'Argol. 
                \item Une fois arrivés, on lancera la préparation du repas et le gonflement des annexes en parallèle. Départ vers la plage à 19h45 pour début de veillée à 20h, lancement du retour sur les bateaux à 21h avant le coucher de soleil.   
            \end{itemize}
        {
        \huge \textbf{Contraintes de la navigation}}\\
            \begin{itemize}
                \item Bien respecter le chenal du Crouesty au départ. 
                \item Si la route dérive vers l'Ouest, il est possible d'arriver au nord du parc à moule de Houat, en ce cas reprendre un cap plus à l'Est pour le laisser largement dans notre Ouest. 
                \item Bien respecter les cardinales Est \textit{Men Groise} et \textit{Er Rouzès} une fois Houat passé. 
                \item Idem pour le danger isolé \textit{La chèvre}. 
                \item Amarrage au port de l'Argol à la tonne, point d'attention sur le parbattage et l'espacement des barres de flèches dans l'axe longitudinale.
                \item L'arrivée se fera à basse mer. Les petits coefficients de marée donnent une hauteur d'eau suffisante pour manoeuvrer sereinement autour des tonnes mais ne pas se priver d'allumer le sondeur s'il fonctionne.   
                \item Contrainte sur la durée de navigation pour pouvoir créer de la vie scout : le départ est contraint par la messe du matin, l'arrivée par la veillée sur la plage d'Hoëdic. 
            \end{itemize}
\end{minipage}%
%%%%%%%%%%%%%%%%
% Page de droite
%%%%%%%%%%%%%%%%
\newgeometry{top=2cm, bottom=2cm, left=2cm, right=2cm}
\begin{minipage}[t][0.9\textheight][t]{0.60\textwidth}
    \color{StrongBlue}    
    \noindent
    \begin{minipage}[t][0.3\textheight][t]{\textwidth}
        \begin{minipage}[t][0.3\textheight][t]{0.65\textwidth}
            \vspace{-8.5em}
            {\huge \textbf{Détails de la navigation}}\\
            Le QR-code situé à droite est cliquable, il permet d'explorer le planning de navigation sur le site datashom. 
            La carte ci-dessous visualise le trajet prévu ainsi que les abris de replis et dangers identifiés. 
        \end{minipage}%
        \hfill
        \begin{minipage}[t][0.3\textheight][t]{0.25\textwidth}       
        \href{https://data.shom.fr/datacontext/fr/1747907917371-909620}{ 
            \includegraphics[width=0.8\textwidth]{qr-code.png}
        }
        \end{minipage}%
    \end{minipage}
    \begin{minipage}[t][0.7\textheight][b]{\textwidth}
        \includegraphics[width=0.97\textwidth]{carte.png}
        % \includegraphics[height=0.7\textheight]{carte.png}
    \end{minipage}
\end{minipage}
\begin{minipage}[t][0.9\textheight][t]{0.40\textwidth}
    \color{StrongBlue} 
    \centering 
    \includegraphics[height=\textheight, width=\textwidth, keepaspectratio]{illustration_verticale.jpg}
\end{minipage}%
\restoregeometry



